\NSPage{103}%
\begin{EESegment}{EE-TGT-P103-001}
Take two waves \NSi{NS-F-P103-001}{w,w'} built from curls, with the same label and with exponents
\NSi{NS-F-P103-002}{\alpha,\alpha'}. The nonzero harmonics of
\NSi{NS-F-P103-003}{(w\mathbin{\cdot}\nabla_*)w'+(w'\mathbin{\cdot}\nabla_*)w} lie in
\NSi{NS-F-P103-004}{\mathsf W_{\alpha+\alpha'-\kappa_s}}, and the zero harmonic lies in the mean
class \NSi{NS-F-P103-005}{\mathsf M_{\alpha+\alpha'-\kappa_s}}. Waves with different labels have
zero product on their closed supports.
\end{EESegment}
\end{lemma}

\begin{proof}
\begin{EESegment}{EE-TGT-P103-002}
When a mean field advects the phase of a wave, the phase derivative costs
\NSi{NS-F-P103-006}{k=O(\varepsilon^{-1/2})}. Radial transport of the amplitude contains a radial
mean factor of order \NSi{NS-F-P103-007}{\mu+1} and one
\NSi{NS-F-P103-008}{D_r}. Axial transport gains one power through
\NSi{NS-F-P103-009}{D_z}. When a wave transports a mean field, the radial cost is at most
\NSi{NS-F-P103-010}{\kappa_s}, and the cylindrical connection terms contain no derivatives.
Each of these terms has at least the claimed mean--wave exponent.
\end{EESegment}

\begin{EESegment}{EE-TGT-P103-003}
Now take two harmonics with the same label. When the transport derivative falls on
\NSi{NS-F-P103-011}{a'\exp(ikm'\Phi)} and the advecting harmonic is
\NSi{NS-F-P103-012}{a\exp(ikm\Phi)}, phase differentiation gives
\NSi{NS-F-P103-013}{ikm'(a\mathbin{\cdot}n_\Phi)a'}. Equation~\eqref{eq:7.39} cancels what looks
like a loss of one half-power. Terms where the derivative falls on the amplitude cost at most
\NSi{NS-F-P103-014}{\kappa_s}, and the connection terms cost nothing. The pulse weights satisfy
\NSi{NS-F-P103-015}{P_v^2\le P_v} and their product has radial weight
\NSi{NS-F-P103-016}{\zeta}. For a nonzero output, this is stronger than the required
\NSi{NS-F-P103-017}{\sqrt\zeta} weight. Lemma~\ref{lem:6.1} gives the support separation for
different labels. Applying Leibniz' rule gives every statement at each fixed order of amplitude
differentiation.
\end{EESegment}
\end{proof}

\begin{EESegment}{EE-TGT-P103-004}
\subsection{Sources and errors kept during the iteration}\label{sec:9.2}
To apply Proposition~\ref{prop:9.1} and Lemma~\ref{lem:9.2} at one stage after another, every new
source has to satisfy the support and derivative assumptions of the pulse inverse in
Proposition~\ref{prop:7.2}. Mean operations can make the auxiliary support larger, so those
assumptions have to be checked after a whole cycle. Equation~\eqref{eq:9.3} keeps the errors that
are already flat at the singularity separate from the sources still waiting to be corrected. Even
so, every nonlinear product uses the full velocity fields. The supported sources also have to stay
inside the fixed enlarged rectangle. Proposition~\ref{prop:9.3} proves all of these facts.
\end{EESegment}

\begin{EESegment}{EE-TGT-P103-005}
Use the source supports \NSi{NS-F-P103-018}{\Omega_\gamma} and the enlarged domains
\NSi{NS-F-P103-019}{\mathcal E_\gamma^+} from \eqref{eq:6.28}. The flat term below contains the
base error \NSi{NS-F-P103-020}{\mathcal E_B} from \eqref{eq:5.41}, the tails caused by pulse
cutoffs, and the cutoff remainders that have not been canceled. The exact mean balances can include
these flat terms, because their definition always uses the full current velocity. Angular means can
fill the whole auxiliary torus. The rectangle and envelope conditions below apply only to the
nonzero harmonics.
\end{EESegment}

\begin{proposition}\label{prop:9.3}
\begin{EESegment}{EE-TGT-P103-006}
Fix \NSi{NS-F-P103-021}{j\in\mathbb N_0}. Suppose
\NSi{NS-F-P103-022}{(u^{[j]},p^{[j]})} is built from the fixed slow base and the primary waves by
a finite sequence of the following operations, the pulse inverse and curl construction, the signed
amplitude map for stresses that do not depend on the auxiliary variables, and the mean maps from
Section~\ref{sec:8}. Reconstruct the pressure through \eqref{eq:8.12} after every update, and use
the full velocity fields in every product. Then the physical residual splits as
\end{EESegment}
\NSFormula{NS-F-P103-023}{display}{9.3}
\begin{equation}
\mathscr R(u^{[j]},p^{[j]})=\mathcal G^{[j]}+\mathcal F^{[j]}.
\tag{9.3}\label{eq:9.3}
\end{equation}
\begin{EESegment}{EE-TGT-P103-007}
In one common chart, set
\NSi{NS-F-P103-024}{\mathcal G_*^{[j]}=Q^{2A+1/2}\mathcal G^{[j]}} and
\NSi{NS-F-P103-025}{\mathcal F_*^{[j]}=Q^{2A+1/2}\mathcal F^{[j]}}. The following three
statements then hold.
\end{EESegment}
\begin{enumerate}[label=(\roman*)]
\item \begin{EESegment}{EE-TGT-P103-008}
\emph{Supported sources.} On the coarsest common auxiliary torus, the nonzero angular harmonics
have the form
\end{EESegment}
\NSFormula{NS-F-P103-026}{display}{9.4}
\begin{equation}
\begin{gathered}
\mathcal G_*^{[j]}-\langle\mathcal G_*^{[j]}\rangle_\theta
 =\sum_{\gamma\in\Gamma}\sum_{m\in\mathscr H_j}
 f_{\gamma,m}^{[j]}e^{ik_\gamma m\Phi_\gamma},\\
\mathscr H_j\subset\mathbb Z\setminus\{0\},\qquad |\mathscr H_j|<\infty.
\end{gathered}
\tag{9.4}\label{eq:9.4}
\end{equation}
\begin{EESegment}{EE-TGT-P103-009}
The sum is locally finite. The set
\NSi{NS-F-P103-027}{\mathscr H_j} depends on the finite construction, but not on the band. The
coefficients satisfy
\end{EESegment}
\NSFormula{NS-F-P103-028}{display}{none}
\[
\operatorname{supp}(f_{\gamma,m}^{[j]}|_{\mathcal E_\gamma})\subset\Omega_\gamma,
\qquad
f_{\gamma,m}^{[j]}\in C^\infty(\mathcal E_\gamma^+;\mathbb C^3),
\]
\begin{EESegment}{EE-TGT-P103-010}
and they extend smoothly by zero across the fixed slow support, the fixed transverse support, and
the enlarged local torus rectangle. Let \NSi{NS-F-P103-029}{\alpha,\mu} be any lower bounds for
the wave and mean exponents obtained by
\end{EESegment}
\NSPage{104}%
\begin{EESegment}{EE-TGT-P104-001}
applying Propositions~\ref{prop:9.1}, \ref{prop:7.2}, \ref{prop:6.6}, and \ref{prop:7.6},
Section~\ref{sec:8}, and Lemma~\ref{lem:9.2} along this finite sequence. Then
\end{EESegment}
\NSFormula{NS-F-P104-001}{display}{none}
\[
f_{\gamma,m}^{[j]}\in \mathsf W_\alpha,
\qquad
\langle\mathcal G_*^{[j]}\rangle_\theta\in \mathsf M_\mu.
\]
\begin{EESegment}{EE-TGT-P104-002}
The mean coefficients extend smoothly by zero outside the active shell. The next pulse inverse is
applied exactly to the coefficients \NSi{NS-F-P104-002}{f_{\gamma,m}^{[j]}} in \eqref{eq:9.4}.
\end{EESegment}

\item \begin{EESegment}{EE-TGT-P104-003}
\emph{Flat remainder.} For every fixed \NSi{NS-F-P104-003}{j,m,N},
\end{EESegment}
\NSFormula{NS-F-P104-004}{display}{9.5}
\begin{equation}
|\mathcal F^{[j]}|_m\le C_{j,m,N}q^N.
\tag{9.5}\label{eq:9.5}
\end{equation}
\begin{EESegment}{EE-TGT-P104-004}
This physical bound is uniform in the band, label, and point.
\end{EESegment}

\item \begin{EESegment}{EE-TGT-P104-005}
\emph{Exact mean equations.} Let
\NSi{NS-F-P104-005}{(\beta^{[j]},v^{[j]},\gamma^{[j]},w^{[j]})} be the full current velocity
correction from \eqref{eq:8.1}, and set
\NSi{NS-F-P104-006}{\mathcal E_{B,*}=Q^{2A+1/2}\mathcal E_B}. Then
\end{EESegment}
\NSFormula{NS-F-P104-007}{display}{9.6}
\begin{equation}
\begin{gathered}
W_{ab}^{[j]}=\langle w_a^{[j]}w_b^{[j]}\rangle_\theta,
\qquad
\mathcal P^{[j]}=\int\langle g_r^{[j]}\rangle_Y\,dR,
\qquad
p_m^{[j]}=\mathsf T_0(g_r^{[j]}-\rho\mathcal P^{[j]}),\\
\langle\mathcal G_*^{[j]}+\mathcal F_*^{[j]}-\mathcal E_{B,*}\rangle_\theta
=\bigl(D_r p_m^{[j]}-g_r^{[j]},E_\theta^{[j]},E_z^{[j]}\bigr),
\end{gathered}
\tag{9.6}\label{eq:9.6}
\end{equation}
\begin{EESegment}{EE-TGT-P104-006}
Here \NSi{NS-F-P104-008}{g_r^{[j]},E_\theta^{[j]},E_z^{[j]}} are the values given by
\eqref{eq:8.3} for this tuple.
\end{EESegment}
\end{enumerate}
\end{proposition}

\begin{proof}
\begin{EESegment}{EE-TGT-P104-007}
First prove (i). Then keep track of the additive errors in (ii) while preserving the exact mean
equations in (iii).
\end{EESegment}

\begin{EESegment}{EE-TGT-P104-008}
\noindent\textbf{Part (i), support and coefficient bounds.} The primary pulses have the stated
supports and weighted bounds at every derivative order. Differentiation preserves both the closed
support and the amplitude-class bounds. Each explicit \NSi{NS-F-P104-009}{D_r} costs
\NSi{NS-F-P104-010}{\kappa_s}. Taking more amplitude derivatives of that estimate can make its
constants and polynomial degrees larger, but it does not charge that exponent loss a second time.
\end{EESegment}

\begin{EESegment}{EE-TGT-P104-009}
In a product of a mean and a wave, the product keeps the wave's slow support, transverse support,
local torus rectangle, and \NSi{NS-F-P104-011}{P_v} factor. Multiplying two waves with one label
adds their integer phase multipliers. Their zero harmonic is exactly their angular mean, because the
angular frequency \NSi{NS-F-P104-012}{kp} is a nonzero integer. Every other output keeps the wave
weight, since \NSi{NS-F-P104-013}{P_v^2\le P_v}. Different labels do not interact. Taking curls,
forming pressure coefficients, and propagating a particular pulse all preserve the phase integer.
Mean operations cannot create a nonzero angular harmonic. Any nonzero harmonic produced later still
contains a wave factor, so it keeps that factor's support and envelope.
\end{EESegment}

\begin{EESegment}{EE-TGT-P104-010}
The radial integration, pressure reconstruction, mean-vector-potential construction, and
auxiliary-time inversion from Section~\ref{sec:8} all preserve the bounds for mean amplitudes. For
the stress correction in Proposition~\ref{prop:7.6}, subtract a supported interior profile with the
same weighted radial moment as the source. The corrected source then has total moment zero. Its
weighted radial integral can therefore be taken forward from the left edge of its support or
backward from the right edge. Near either edge, let \NSi{NS-F-P104-014}{s} be the matching
logarithmic distance from \eqref{eq:6.23}. There, the weight \NSi{NS-F-P104-015}{\zeta} is
\NSi{NS-F-P104-016}{e^{-a/s^2}} times a smooth positive factor with fixed positive upper and lower
bounds, where \NSi{NS-F-P104-017}{a>0} is fixed. For every fixed
\NSi{NS-F-P104-018}{M\ge0} and every integer \NSi{NS-F-P104-019}{k\ge0}, sufficiently small
positive \NSi{NS-F-P104-020}{\delta,s} satisfy
\end{EESegment}
\NSFormula{NS-F-P104-021}{display}{none}
\[
\int_0^\delta s^{-M}e^{-a/s^2}\,ds
 \le C_{M,a}\delta^{3-M}e^{-a/\delta^2},
\qquad
\left|\frac{d^k}{ds^k}e^{-a/s^2}\right|
 \le C_{k,a}s^{-3k}e^{-a/s^2}.
\]
\NSPage{105}%
\begin{EESegment}{EE-TGT-P105-001}
The substitution \NSi{NS-F-P105-001}{v=a/s^2} gives the first bound, and direct differentiation
gives the second. On the normalized shell, the radial weights and the radial Jacobian are bounded.
Differentiating the radial integral therefore leaves a bound of the form
\NSi{NS-F-P105-002}{\zeta\delta^{-M'}}, where \NSi{NS-F-P105-003}{M'} can depend on the
derivative. In the interior, \NSi{NS-F-P105-004}{\zeta} has a positive lower bound. A source of
order \NSi{NS-F-P105-005}{\alpha} therefore gives a stress correction
\NSi{NS-F-P105-006}{\Sigma\in \mathsf M_\alpha}.
\end{EESegment}

\begin{EESegment}{EE-TGT-P105-002}
For the amplitude division, use the same normalization as in \eqref{eq:7.31}, namely
\NSi{NS-F-P105-007}{d_\Sigma=\mathsf H^{-1}(\Sigma/\varepsilon)} and
\NSi{NS-F-P105-008}{\delta a_\sigma=(d_\Sigma)_\sigma/(2a_\sigma)}. The inverse-matrix bounds and
\eqref{eq:7.25} give the following estimates for every fixed coefficient multi-index
\NSi{NS-F-P105-009}{I}.
\end{EESegment}
\NSFormula{NS-F-P105-010}{display}{none}
\[
\begin{aligned}
|D^I(d_\Sigma)_\sigma|
 &\le C_I\varepsilon^{\alpha-1}S_*^{b_I}\zeta\delta^{-M_I},\\
|D^I(a_\sigma^{-1})|
 &\le C_I S_*^{b_I}\zeta^{-1/2}\delta^{-M_I},\\
|D^I\delta a_\sigma|
 &\le C_I\varepsilon^{\alpha-1}S_*^{b'_I}\sqrt\zeta\,\delta^{-M'_I},\\
|D^I(\sqrt\varepsilon\,\delta a_\sigma b_\sigma)|
 &\le C_I\varepsilon^{\alpha-1/2}S_*^{b''_I}\sqrt\zeta\,\delta^{-M''_I}P_v.
\end{aligned}
\]
\begin{EESegment}{EE-TGT-P105-003}
This gives \eqref{eq:7.33} again. The last line also uses the fundamental pulse bounds. After the slow
cutoff, it is the \NSi{NS-F-P105-011}{\mathsf W_{\alpha-1/2}} estimate from
Proposition~\ref{prop:7.6}. The remaining square-root weight at the edge makes the extension by zero
smooth.
\end{EESegment}

\begin{EESegment}{EE-TGT-P105-004}
Each new source now satisfies the assumptions of Proposition~\ref{prop:7.2}. Along its prescribed
path, the slow and transverse variables are fixed. If the data are zero all along one of these
paths, the solution is zero there as well. The final cutoff puts the support back inside the fixed
enlarged local torus rectangle. A source with a smooth extension gives a solution with a smooth
extension. This proves the support condition needed by the inverse. There is no need to divide by
the original cutoff.
\end{EESegment}

\begin{EESegment}{EE-TGT-P105-005}
The set of local bands stays fixed. Radial paths keep \NSi{NS-F-P105-012}{(z,t)} fixed, and
\NSi{NS-F-P105-013}{q} depends only on \NSi{NS-F-P105-014}{(z,t)}. Temporal inversion and pulse
propagation keep all slow variables fixed. The common-torus construction from Lemma~\ref{lem:6.2}
therefore applies at every step of the finite construction. A radial mean can collect a number of
labels that grows polynomially in \NSi{NS-F-P105-015}{S_*}, while only a bounded number of wave
labels overlap at any one point. This radial collection changes only the polynomial factors in
\NSi{NS-F-P105-016}{S_*}. It does not change the common chart or the
\NSi{NS-F-P105-017}{\varepsilon} exponent.
\end{EESegment}

\begin{EESegment}{EE-TGT-P105-006}
\noindent\textbf{Parts (ii) and (iii), flat errors and exact means.} When an actual increment is
\NSi{NS-F-P105-018}{(w,\pi)}, the residual changes by
\end{EESegment}
\NSFormula{NS-F-P105-019}{display}{none}
\[
\begin{aligned}
\mathscr R(u+w,p+\pi)={}&\mathscr R(u,p)+\partial_t w-\Delta w+\nabla\pi\\
&+(u\mathbin{\cdot}\nabla)w+(w\mathbin{\cdot}\nabla)u+(w\mathbin{\cdot}\nabla)w.
\end{aligned}
\]
\begin{EESegment}{EE-TGT-P105-007}
The old \NSi{NS-F-P105-020}{\mathcal F} is added directly. Every product in the display uses the
actual fields, including the curls with their cutoffs and the remainders from the modified radial
integral with compact support. Put every newly created flat term into
\NSi{NS-F-P105-021}{\mathcal F}. At any fixed stage, there are only finitely many kinds of these
terms. The Gaussian bound in Proposition~\ref{prop:9.1} is uniform over the labels.
Lemma~\ref{lem:8.2} bounds each cutoff remainder by any chosen power of
\NSi{NS-F-P105-022}{q}, at the cost of using finitely many more source derivatives. Summing a
number of labels that grows polynomially does not spoil those bounds. Keep the Gaussian cutoff tails
with nonzero harmonics in the additive residual, and do not feed them into later forward pulse
solves. The mean residuals \NSi{NS-F-P105-023}{E_\theta,E_z} are the exact conservative balances
and may contain flat cutoff remainders, which a fixed mean inverse may cancel. If a contribution
is canceled, remove it from \NSi{NS-F-P105-024}{\mathcal F}. The cutoff-remainder estimates put the
remainder in \NSi{NS-F-P105-025}{\mathsf M_\alpha} for every
\NSi{NS-F-P105-026}{\alpha}. This update keeps both the stated classes and the exact splitting of
the residual. It proves \eqref{eq:9.5} at every finite stage, which proves (ii). Apply
Proposition~\ref{prop:8.1} to the full current tuple and use \eqref{eq:8.12}. This gives
\eqref{eq:9.6} and proves (iii).
\end{EESegment}
\end{proof}

\begin{EESegment}{EE-TGT-P105-008}
In the realization step, the finite corrections are summed to build the infinite local field. Its
residual is controlled by comparing that field with one fixed finite stage. The flat residuals
from the separate stages are not added together.
\end{EESegment}

\NSPage{106}%
\begin{EESegment}{EE-TGT-P106-001}
\subsection{Starting the iteration and running a full correction cycle}\label{sec:9.3}
The primary pulses, followed by the first mean corrections, give the first residual bounds in
Proposition~\ref{prop:9.5}. Proposition~\ref{prop:9.6} then shows that a full cycle improves every
one of those bounds while keeping the support and moment conditions. The induction statement below
combines the residual orders in \eqref{eq:9.8} with the bounds on the total velocity correction in
\eqref{eq:9.9}. Those total bounds control how the existing correction interacts with every later
increment.
\end{EESegment}

\begin{EESegment}{EE-TGT-P106-002}
Throughout the construction, keep the base \NSi{NS-F-P106-001}{(b,V,G)}, the primary amplitudes,
and every inverse operator fixed. Let
\NSi{NS-F-P106-002}{w_0^{\mathrm{tan}}=W_0^{\mathrm{as}}} be the normalized representative of the
assembled primary transverse field from \eqref{eq:7.35}. For each finite construction, write its
normalized velocity and pressure as
\end{EESegment}
\NSFormula{NS-F-P106-003}{display}{9.7}
\begin{equation}
u_*=(b+\beta,V+v,G+\gamma)+w,
\qquad
p_*=p_{B,*}+p_m+p_w,
\qquad
\langle w\rangle_\theta=\langle p_w\rangle_\theta=0.
\tag{9.7}\label{eq:9.7}
\end{equation}
\begin{EESegment}{EE-TGT-P106-003}
Here \NSi{NS-F-P106-004}{p_{B,*}} is the pressure of the base,
\NSi{NS-F-P106-005}{p_w} is the sum of the pressure harmonics supplied by the amplitude equations,
and \NSi{NS-F-P106-006}{p_m} is reconstructed from the current
\NSi{NS-F-P106-007}{g_r} through \eqref{eq:8.12}. Equations~\eqref{eq:8.3}, \eqref{eq:8.12}, and
\eqref{eq:8.15} then determine
\NSi{NS-F-P106-008}{E_\theta,E_z,\mathcal P,\mathcal J_\theta,\mathcal J_z}.
\end{EESegment}

\begin{definition}[Finite correction state]\label{def:9.4}
\begin{EESegment}{EE-TGT-P106-004}
For \NSi{NS-F-P106-009}{j\in\mathbb N\cup\{0\}}, a state at stage
\NSi{NS-F-P106-010}{j} consists of real fields of the form \eqref{eq:9.7} with the properties
listed next. The wave \NSi{NS-F-P106-011}{w} is a sum of curls of supported wave potentials with
the fixed phases for each label. The mean correction
\NSi{NS-F-P106-012}{(\beta,v,\gamma)} is the sum of the curl of an azimuthal potential and an
azimuthal field added directly. Both velocity corrections are therefore exactly divergence-free.
The source supports, finite sets of harmonics, smooth extensions by zero, and physical residual
splitting are the ones in Proposition~\ref{prop:9.3}.
\end{EESegment}

\begin{EESegment}{EE-TGT-P106-005}
The normalized nonzero harmonics \NSi{NS-F-P106-013}{\mathcal G_{\mathrm{wave}}^{[j]}} of
\NSi{NS-F-P106-014}{\mathcal G_*^{[j]}}, the exact tangential mean residuals, and the three
integral defects obey
\end{EESegment}
\NSFormula{NS-F-P106-015}{display}{9.8}
\begin{equation}
\begin{gathered}
\mathcal G_{\mathrm{wave}}^{[j]}\in \mathsf W_{B_j},
\qquad E_\theta,E_z\in \mathsf M_{C_j^*},
\qquad (\mathcal P,\mathcal J_\theta,\mathcal J_z)\in \mathsf S_{C_j^*},\\
B_j=\frac12+\sigma_j,
\qquad C_j^*=1+\sigma_j,
\qquad \sigma_j=\frac15+\frac{j}{10},
\end{gathered}
\tag{9.8}\label{eq:9.8}
\end{equation}
\begin{EESegment}{EE-TGT-P106-006}
and the total corrections to velocity and pressure satisfy
\end{EESegment}
\NSFormula{NS-F-P106-016}{display}{9.9}
\begin{equation}
w\in \mathsf W_{1/2},
\qquad w-w_0^{\mathrm{tan}}\in \mathsf W_{0.68},
\qquad v,\gamma,p_m\in \mathsf M_{0.9},
\qquad \beta\in \mathsf M_{1.9}.
\tag{9.9}\label{eq:9.9}
\end{equation}
\begin{EESegment}{EE-TGT-P106-007}
Each wave statement applies to every grouped nonzero harmonic for each label. The two normalized
moments for angular momentum and axial flux stay exactly equal to zero,
\end{EESegment}
\NSFormula{NS-F-P106-017}{display}{9.10}
\begin{equation}
\int_0^\infty R^2\langle v\rangle_Y\,dR=0,
\qquad
\int_0^\infty R\langle\gamma\rangle_Y\,dR=0.
\tag{9.10}\label{eq:9.10}
\end{equation}
\begin{EESegment}{EE-TGT-P106-008}
Reconstructing the pressure is part of the state. It follows that the radial residual of every such
state is \NSi{NS-F-P106-018}{-\rho\mathcal P} plus the flat cutoff remainder from
Proposition~\ref{prop:8.3}.
\end{EESegment}
\end{definition}

\begin{proposition}\label{prop:9.5}
\begin{EESegment}{EE-TGT-P106-009}
Start with the fixed base and the curls of the primary potentials. Apply the temporal mean update
\eqref{eq:8.20}, then apply the five-equation correction \eqref{eq:8.25}. Reconstruct the pressure
through \eqref{eq:8.12} both before and after each update. The resulting fields form a stage
\NSi{NS-F-P106-019}{0} state in the sense of Definition~\ref{def:9.4}. In particular, their
residual orders are \NSi{NS-F-P106-020}{B_0=0.7} and
\NSi{NS-F-P106-021}{C_0^*=1.2}.
\end{EESegment}
\end{proposition}

\begin{proof}
\begin{EESegment}{EE-TGT-P106-010}
First estimate the primary waves. Then remove the means that depend on the auxiliary torus and the
three moment defects. This first stage uses the faster decay of the primary residual after taking
its auxiliary average. The primary transverse amplitudes have exponent
\NSi{NS-F-P106-022}{1/2}. The supported parts of their linear residual have exponent
\NSi{NS-F-P106-023}{1-3\kappa_s}, and the nonlinear residuals of the waves have
\end{EESegment}
\NSPage{107}%
\begin{EESegment}{EE-TGT-P107-001}
exponent \NSi{NS-F-P107-001}{1-\kappa_s}. Before any mean update, the exact mean balances in
\eqref{eq:8.3} give \NSi{NS-F-P107-002}{g_r,p_m,E_\theta,E_z\in \mathsf M_{1-\kappa_s}} and put
the defects in \NSi{NS-F-P107-003}{\mathsf S_{1-\kappa_s}}.
\end{EESegment}

\begin{EESegment}{EE-TGT-P107-002}
Now take the auxiliary means. The primary covariance from Proposition~\ref{prop:7.5}, assembled by
\eqref{eq:7.35}, cancels the order-zero auxiliary radial--tangential stress tensor exactly. This
also includes derivatives of the slow partition, because the squares of the partition functions add
up to one. To write out this cancellation, let
\NSi{NS-F-P107-004}{\Sigma_a^{(0)}} be the leading tensor and let
\NSi{NS-F-P107-005}{\Sigma_a} be the full tensor in \eqref{eq:8.3}. Before adding the mean
velocities, those equations become
\end{EESegment}
\NSFormula{NS-F-P107-006}{display}{9.11}
\begin{equation}
\begin{aligned}
\langle E_\theta\rangle_Y
={}&(\partial_R+2/R)(\langle W_{r\theta}\rangle_Y-\Sigma_\theta^{(0)})
 -(\partial_R+2/R)(\Sigma_\theta-\Sigma_\theta^{(0)})
 +\varepsilon\partial_Z\langle W_{z\theta}\rangle_Y,\\
\langle E_z\rangle_Y
={}&(\partial_R+1/R)(\langle W_{rz}\rangle_Y-\Sigma_z^{(0)})
 -(\partial_R+1/R)(\Sigma_z-\Sigma_z^{(0)})
 +\varepsilon\partial_Z\langle W_{zz}+p_m\rangle_Y.
\end{aligned}
\tag{9.11}\label{eq:9.11}
\end{equation}
\begin{EESegment}{EE-TGT-P107-003}
Each first difference contains at least one curl correction. Before taking the divergence, the tensor
cross term between a curl correction and a primary transverse amplitude lies in
\NSi{NS-F-P107-007}{\mathsf M_{3/2-\kappa_s}}. The axial fluxes, the axial pressure term, and the
higher-order auxiliary stress tensor all have residual exponent at least
\NSi{NS-F-P107-008}{2-\kappa_s}. It follows that
\NSi{NS-F-P107-009}{\langle E_\theta\rangle_Y,\langle E_z\rangle_Y\in \mathsf M_{1.49}}.
\end{EESegment}

\begin{EESegment}{EE-TGT-P107-004}
Set \NSi{NS-F-P107-010}{H_0=1-\kappa_s} and apply the temporal update \eqref{eq:8.20}. Its
tangential increment lies in \NSi{NS-F-P107-011}{\mathsf M_{H_0}}, and the radial velocity it
creates lies in \NSi{NS-F-P107-012}{\mathsf M_{H_0+1}}. The pressure change has exponent
\NSi{NS-F-P107-013}{H_0}, because \NSi{NS-F-P107-014}{2V\Delta v/R} occurs in
\NSi{NS-F-P107-015}{\Delta g_r}. In either tangential residual, the contribution from this pressure
change contains an axial derivative and gains the matching extra power. Once the fast-time term has
been canceled, a slow-time derivative supplies one more power of
\NSi{NS-F-P107-016}{\varepsilon}. Each radial linear flux contains either
\NSi{NS-F-P107-017}{b=O(\varepsilon)} in every fixed-order amplitude derivative on the shell, or
the radial velocity increment that was just created. Axial fluxes gain one power, while viscosity
gains \NSi{NS-F-P107-018}{1-2\kappa_s}. Every radial quadratic mean flux also contains a radial
mean. All these changes to the residual have exponent at least
\NSi{NS-F-P107-019}{H_0+1-2\kappa_s}. Interaction of the mean correction with
\NSi{NS-F-P107-020}{w\in \mathsf W_{1/2}} has wave exponent at least
\NSi{NS-F-P107-021}{H_0}. After pressure is recomputed, the defects still lie in
\NSi{NS-F-P107-022}{\mathsf S_{H_0}}.
\end{EESegment}

\begin{EESegment}{EE-TGT-P107-005}
Next apply the five-equation correction map \eqref{eq:8.25} at order
\NSi{NS-F-P107-023}{H_0}, and recompute the pressure. Its slow increments have no fast-time term,
so the tangential and wave estimates from the preceding paragraph still hold. In the changes to the
defects, \NSi{NS-F-P107-024}{2V\Delta v/R} and its two contributions to the weighted flux moments
are exactly the linear terms in
\NSi{NS-F-P107-025}{\Delta\mathcal P,\Delta\mathcal J_\theta,\Delta\mathcal J_z} that this system
cancels. Every term left contains either a slow-time derivative of the radial velocity increment, a
radial flux with \NSi{NS-F-P107-026}{b} or another radial mean, an axial derivative, radial
viscosity, or one more tangential mean in
\NSi{NS-F-P107-027}{v^2,\gamma v,\gamma^2}. Each has a gain greater than
\NSi{NS-F-P107-028}{0.8} over \NSi{NS-F-P107-029}{H_0}. The remaining defects therefore have
exponent greater than \NSi{NS-F-P107-030}{1.8-\kappa_s}, and the two tangential means keep
exponent \NSi{NS-F-P107-031}{1.49}. These estimates give \eqref{eq:9.8} with
\NSi{NS-F-P107-032}{B_0=0.7} and \NSi{NS-F-P107-033}{C_0^*=1.2}.
\end{EESegment}

\begin{EESegment}{EE-TGT-P107-006}
The primary curl correction has exponent
\NSi{NS-F-P107-034}{1-\kappa_s>0.68}. Every mean increment and pressure change has exponent
\NSi{NS-F-P107-035}{H_0>0.9}, and its radial exponent is
\NSi{NS-F-P107-036}{H_0+1>1.9}. This proves \eqref{eq:9.9}. The prescribed temporal increments
have auxiliary mean zero. The first two homogeneous equations in the correction system enforce
\eqref{eq:9.10}, and reconstructing the mean field preserves the prescribed auxiliary mean
exactly.
\end{EESegment}
\end{proof}

\begin{EESegment}{EE-TGT-P107-007}
The initial stage now has the required bounds. Repeat the same four corrections, and compare every
new error with the bound needed at the next stage. When estimating how the new increments interact
with the existing fields, the total bounds in \eqref{eq:9.9} remain available.
\end{EESegment}

\begin{proposition}\label{prop:9.6}
\begin{EESegment}{EE-TGT-P107-008}
Let \NSi{NS-F-P107-037}{(u^{[j]},p^{[j]})} be a stage
\NSi{NS-F-P107-038}{j} state as in Definition~\ref{def:9.4}. Using the same base, primary
amplitudes, supports, and inverse operators, it is possible to construct a stage
\NSi{NS-F-P107-039}{(u^{[j+1]},p^{[j+1]})} state whose residual orders are
\end{EESegment}
\NSFormula{NS-F-P107-040}{display}{none}
\[
B_{j+1}=B_j+\frac1{10},
\qquad
C_{j+1}^*=C_j^*+\frac1{10}.
\]
\begin{EESegment}{EE-TGT-P107-009}
Construct it by applying the following corrections in this order.
\end{EESegment}
\begin{enumerate}[label=(\roman*)]
\NSPage{108}%
\item \begin{EESegment}{EE-TGT-P108-001}
Solve the inhomogeneous amplitude equation for each supported nonzero harmonic.
\end{EESegment}
\item \begin{EESegment}{EE-TGT-P108-002}
Change the wave amplitudes so that they correct the tangential residual after taking its auxiliary
average.
\end{EESegment}
\item \begin{EESegment}{EE-TGT-P108-003}
Apply the temporal mean inverse to the part whose auxiliary average is zero.
\end{EESegment}
\item \begin{EESegment}{EE-TGT-P108-004}
Apply the five-equation correction to the three current defects.
\end{EESegment}
\end{enumerate}
\begin{EESegment}{EE-TGT-P108-005}
Reconstruct the pressure after every correction. Build the wave increments and the meridional
velocity increments, meaning the radial and axial velocity increments, from their potentials, and
add the azimuthal mean increments directly. The cumulative bounds
\eqref{eq:9.9} and the exact moments \eqref{eq:9.10} are preserved.
\end{EESegment}
\end{proposition}

\begin{proof}
\begin{EESegment}{EE-TGT-P108-006}
The four steps below carry out (i)--(iv) in that order. For this cycle, shorten the notation to
\NSi{NS-F-P108-001}{B=\frac12+\sigma_j} and
\NSi{NS-F-P108-002}{C_*=B+\frac12=1+\sigma_j}. Whenever
\NSi{NS-F-P108-003}{g_r} or \NSi{NS-F-P108-004}{p_m} appears below, an improved order refers to
the change in that quantity. The total corrections still satisfy \eqref{eq:9.9}. Within each step,
\NSi{NS-F-P108-005}{w,\beta,v,\gamma} are the fields at the start of that step. The increment
replaces them by
\end{EESegment}
\NSFormula{NS-F-P108-006}{display}{none}
\[
w+\Delta w,
\qquad
(\beta+\Delta\beta,v+\Delta v,\gamma+\Delta\gamma).
\]
\begin{EESegment}{EE-TGT-P108-007}
Use these full updated fields in \eqref{eq:8.3} to compute
\NSi{NS-F-P108-007}{g_{r,\mathrm{new}}}, and set
\end{EESegment}
\NSFormula{NS-F-P108-008}{display}{none}
\[
\mathcal P_{\mathrm{new}}=\int\langle g_{r,\mathrm{new}}\rangle_Y\,dR,
\qquad
p_{m,\mathrm{new}}=\mathsf T_0(g_{r,\mathrm{new}}-\rho\mathcal P_{\mathrm{new}}).
\]
\begin{EESegment}{EE-TGT-P108-008}
Compute every residual and defect in the next step from this updated state.
\end{EESegment}

\begin{EESegment}{EE-TGT-P108-009}
\noindent\textbf{Step 1, cancel the supported harmonics.} Apply
Proposition~\ref{prop:7.2} to each grouped supported source. Use zero initial data where the
prescribed path enters the common auxiliary torus. For a source coefficient
\NSi{NS-F-P108-009}{f_m}, the normalized harmonic increments are
\end{EESegment}
\NSFormula{NS-F-P108-010}{display}{none}
\[
\Delta w_m=\operatorname{curl}_*\!\left(
 \frac{i n_\Phi\times(\psi t_m)}{km|n_\Phi|^2}e^{ikm\Phi}\right),
\qquad
\Delta p_{w,m}=\psi\pi_m e^{ikm\Phi},
\qquad
\mathscr L_m(t_m,\pi_m)=-f_m.
\]
\begin{EESegment}{EE-TGT-P108-010}
Put these fields together with their physical rescalings over every label and harmonic, keeping each
conjugate pair. The velocity increment lies in \NSi{NS-F-P108-011}{\mathsf W_B}, and its pressure
lies in \NSi{NS-F-P108-012}{\mathsf W_{B+1/2}}. The new errors with nonzero harmonics have these
exponents.
\end{EESegment}
\NSFormula{NS-F-P108-013}{display}{none}
\[
\begin{array}{r|l}
\text{linear error} & B+1/2-3\kappa_s\\
\text{cross with old exact waves} & B+1/2-\kappa_s\\
\text{particular-correction self-interaction} & 2B-\kappa_s\\
\text{cross with total mean correction} & B+0.4
\end{array}
\]
\begin{EESegment}{EE-TGT-P108-011}
Proposition~\ref{prop:9.1} gives the first row, and Lemma~\ref{lem:9.2} gives the other three. The
change in mean covariance caused by these increments has exponent at least
\NSi{NS-F-P108-014}{B+1/2=C_*}. Taking a radial divergence costs
\NSi{NS-F-P108-015}{\kappa_s}. The exact balances and the linear pressure map now give
\end{EESegment}
\NSFormula{NS-F-P108-016}{display}{none}
\[
E_\theta,E_z,\Delta g_r,\Delta p_m\in \mathsf M_{C_*-\kappa_s},
\qquad
(\mathcal P,\mathcal J_\theta,\mathcal J_z)\in \mathsf S_{C_*-\kappa_s}.
\]
\begin{EESegment}{EE-TGT-P108-012}
The angular-momentum and axial-flux constraints remain unchanged. Together with the exact integrated
identities in \eqref{eq:8.16}, this improves the weighted moments by one power,
\end{EESegment}
\NSFormula{NS-F-P108-017}{display}{9.12}
\begin{equation}
\int R^2\langle E_\theta\rangle_Y\,dR,
\qquad
\int R\langle E_z\rangle_Y\,dR
\in \mathsf S_{C_*+1-\kappa_s}.
\tag{9.12}\label{eq:9.12}
\end{equation}
\begin{EESegment}{EE-TGT-P108-013}
The axial identity contains \NSi{NS-F-P108-018}{c_\rho\mathcal P} inside its axial derivative.
This is the term that accounts for the normalized correction
\NSi{NS-F-P108-019}{-\rho\mathcal P} left in the radial equation after the pressure has been
reconstructed.
\end{EESegment}

\begin{EESegment}{EE-TGT-P108-014}
\noindent\textbf{Step 2, correct the auxiliary-averaged residual by changing the stress.} Once its
total weighted radial moments have been removed, a stress with compact support can cancel the
tangential residual after taking its auxiliary average. Make this correction in physical
coordinates, so the stress agrees between charts. Use the residual components
\NSi{NS-F-P108-020}{E_\theta,E_z} from \eqref{eq:8.3}, supported in the prescribed annulus, and
keep the background residual separate with its flatness
\end{EESegment}
